% Additional proved subsection, to follow moderate_reflection.tex.
% It uses that section's notation and does not enlarge the half-term theorem.
\subsection{A small proportional range and a second coefficient transition}
\label{moderate:sec:proportional}

The preceding circle argument also gives a nonzero proportional range
for the late coefficients themselves. This strengthens the Gaussian
formula without changing the stated range of the half-term remainder
law.

\begin{theorem}[Uniform small proportional saddle]
\label{moderate:thm:proportional}
There exists $c_0\in(0,1)$ with the following properties. For
$0\le c\le c_0$ there is a unique real-analytic continuation
$\tau_c$ of $\tau_0=\tau$ solving
\begin{equation}\label{moderate:eq:moving-saddle}
 D\log\phi(\tau_c)+cD\log a(\tau_c)=1.
\end{equation}
Its curvature
\[
 V(c)=D^2\log\phi(\tau_c)+cD^2\log a(\tau_c)
\]
is positive. Define
\begin{equation}\label{moderate:eq:rate-prefactor}
 S(c)=\log\phi(\tau_c)+c\log a(\tau_c)-\log\tau_c,
 \qquad C(c)=\frac{\tau_c}{\sqrt{2\pi V(c)}}.
\end{equation}
Uniformly for integers $0\le m\le c_0 k$, with $c=m/k$,
\begin{equation}\label{moderate:eq:proportional}
 A_{k,m}=(-1)^{k+m-1}C(c)k^{-3/2}e^{kS(c)}
                         \bigl(1+O(k^{-1})\bigr).
\end{equation}
There is a uniform expansion to every fixed order in integer powers of
$k^{-1}$, with real-analytic coefficient functions of $c$.
The constant $c_0$ is asserted to exist; no numerical lower bound for it
or maximality assertion is included.
\end{theorem}

\begin{proof}
Use the logarithmic radius $x=\log u$ and put
$f(x)=\log\phi(e^x)-x$, $g(x)=\log a(e^x)$ locally near
$x_0=\log\tau$. Then $f'(x_0)=0$ and $f''(x_0)=v>0$.
The analytic implicit-function theorem applied to
$f'(x)+cg'(x)=0$ supplies $x(c)=\log\tau_c$ near zero.
After shrinking the interval, $\tau_c$ lies in a fixed compact
$I\subset(0,1)$ where $a$ is positive, and $V(c)>v/2$.
Uniqueness here means uniqueness of this local continuation; it does
not classify other real or complex saddles.

Choose a small fixed angular neighborhood $|t|<\epsilon$ on which
$a(ue^{it})$ is nonzero for every $u\in I$. Positivity of the
Taylor coefficients of $\phi$ gives a strict uniform bound
\[
 q:=\sup_{u\in I,\ \epsilon\le|t|\le\pi}
         \frac{|\phi(ue^{it})|}{\phi(u)}<1.
\]
Also let
\[
 H=\max\left\{1,\sup_{u\in I,\ |t|\le\pi}
                   \frac{|a(ue^{it})|}{a(u)}\right\}<\infty.
\]
Further decrease $c_0$ so that $qH^{c_0}<1$. In the Cauchy integral
on $|u|=\tau_c$, the outer arc divided by the positive saddle
amplitude is at most $q^kH^m\le(qH^{c_0})^k$.
This is the global estimate needed for proportional $m$: no logarithm
or fractional power of $a$ on the full circle is assumed.

On the small arc the normalized local logarithmic phase is
\[
 f(x(c)+it)-f(x(c))
     +c\bigl(g(x(c)+it)-g(x(c))\bigr).
\]
Its linear term vanishes by \eqref{moderate:eq:moving-saddle}; its
quadratic term is $-V(c)t^2/2$. Analyticity, reality on the real
radius, and compactness give a uniform Gaussian modulus bound for
a sufficiently small fixed $\epsilon$. Choose this $\epsilon$ first
and then make the preceding final reduction of $c_0$, so the outer-arc
estimate uses the same angular neighborhood. Taylor expansion with $t=z/\sqrt k$ and
amplitude $e^{it}$ now gives
\[
 \int e^{k\{f(x(c)+it)-f(x(c))
                +c[g(x(c)+it)-g(x(c))]\}}e^{it}\,dt
 =\sqrt{\frac{2\pi}{kV(c)}}\bigl(1+O(k^{-1})\bigr).
\]
The discarded arcs are exponentially small. Odd terms integrate to
zero, so only integer powers of $k^{-1}$ occur. Analyticity and
uniform derivative bounds also prove the asserted all-orders
expansion. Restoring the Cauchy prefactor
$(-1)^{k+m-1}\tau_c e^{kS(c)}/(2\pi k)$ proves the theorem.
\end{proof}

\begin{corollary}[Explicit rate expansion and cubic transition]
\label{moderate:cor:cubic-transition}
Put
\begin{align*}
 d_1&=-\frac{r_1+r_2/2}{v}+\frac{\kappa_3r_1}{2v^2},\\
 d_3&=\frac{r_2r_1^2}{2v^2}-\frac{\kappa_3r_1^3}{6v^3}.
\end{align*}
Then
\begin{align}\label{moderate:eq:analytic-rate}
 S(c)&=\Lambda+c\log b-\delta_\Gamma c^2+d_3c^3+O(c^4),\notag\\
 C(c)&=C\bigl(1+d_1c+O(c^2)\bigr).
\end{align}
Consequently the Gaussian-resummed leading expression
\begin{equation}\label{moderate:eq:extended-Gaussian}
 A_{k,m}\sim(-1)^{k+m-1}Cb^m\rho^{-k}k^{-3/2}
                     e^{-\delta_\Gamma m^2/k}
\end{equation}
is uniform for $m=o(k^{2/3})$: precisely, it holds uniformly for
$0\le m\le\varepsilon_k k^{2/3}$ whenever $\varepsilon_k\to0$.
For each fixed $L<\infty$, uniformly for $0\le m\le Lk^{2/3}$,
\begin{equation}\label{moderate:eq:cubic-uniform}
 A_{k,m}=(-1)^{k+m-1}Cb^m\rho^{-k}k^{-3/2}
  \exp\left(-\frac{\delta_\Gamma m^2}{k}
               +\frac{d_3m^3}{k^2}\right)
           \bigl(1+O_L(k^{-1/3})\bigr).
\end{equation}
In particular, if $m/k^{2/3}\to\chi\ge0$, the ratio to the right-hand
leading expression in \eqref{moderate:eq:extended-Gaussian} tends to
$e^{d_3\chi^3}$. The numerical value is
$d_3=-0.00545592970200504\ldots$; the exact defining formula is used
in the theorem.
\end{corollary}

\begin{proof}
Differentiating the saddle equation at zero gives
$x'(0)=-r_1/v$. Since $S(c)=f(x(c))+cg(x(c))$ and
$f'(x(c))+cg'(x(c))=0$, expansion through degree three gives
\[
 S(c)=\Lambda+c\log b-\frac{r_1^2}{2v}c^2+
       \left(\frac{r_2r_1^2}{2v^2}
             -\frac{\kappa_3r_1^3}{6v^3}\right)c^3+O(c^4).
\]
For the prefactor,
\[
 (\log C)'(0)=x'(0)-\frac{\kappa_3x'(0)+r_2}{2v}=d_1.
\]
This proves \eqref{moderate:eq:analytic-rate}. Upon setting $c=m/k$,
Theorem~\ref{moderate:thm:proportional} becomes
\begin{align*}
 A_{k,m}={}&(-1)^{k+m-1}Cb^m\rho^{-k}k^{-3/2}
 e^{-\delta_\Gamma m^2/k}\,
 \exp\left(\frac{d_3m^3}{k^2}+O\left(\frac{m^4}{k^3}\right)\right)\\
 &\hspace{4em}\times\left(1+\frac{d_1m}{k}
                  +O\left(\frac{m^2}{k^2}+\frac1k\right)\right).
\end{align*}
For $m=o(k^{2/3})$ all corrections on the last two factors tend to
one uniformly. For $m\le Lk^{2/3}$, the quartic rate error and the
first prefactor correction are $O_L(k^{-1/3})$, which gives
\eqref{moderate:eq:cubic-uniform} and its limiting assertion.
Finally, when $m=\xi\sqrt k$, the prefactor and cubic rate contribute
$(d_1\xi+d_3\xi^3)/\sqrt k$. This is exactly $P_1(\xi)$ in
\eqref{moderate:eq:P1}, independently matching the square-root
calculation.
\end{proof}

\begin{proposition}[Strict cubic damping]
\label{moderate:prop:cubic-negative}
The exact cubic rate coefficient satisfies $d_3<0$.
Thus the transition factor $e^{d_3\chi^3}$ is strictly below one
for every $\chi>0$.
\end{proposition}
\begin{proof}
We already know $\tau>1/2$ and $r_1<0$. The positive log-gamma
series also gives
\[
 D\log\phi(3/5)>
 \frac12\frac35+\frac{(3/5)^2}{1-3/5}
 =\frac65>1,
\]
so $\tau<3/5$. Put $P=\psi(1+\tau)$ and
$Q=\psi'(1+\tau)$. The trigamma series and its
derivatives~\cite{moderate:DLMF} show that $\psi$ is increasing
and concave and that
\[
 Q=\sum_{n\ge0}\frac1{(1+\tau+n)^2}>
       \frac1{1+\tau}.
\]
Using $\gamma>1/2$, $\log2>2/3$, and $\pi^2<10$ in the
classical values at $3/2$ gives
\[
 0<P<\psi(3/2)+\frac1{10}\psi'(3/2)
       <\frac16+\frac1{10}=\frac4{15},
 \qquad
 \tau Q>\frac{\tau}{1+\tau}>\frac13.
\]
In particular $P<\tau Q$.

Write $T=\psi'(1-\tau)>0$. Termwise comparison in the
polygamma series yields
$-\psi''(1-\tau)\le2T/(1-\tau)$. At the saddle,
\[
 v=1+\tau^2T,\qquad
 \kappa_3=1+3\tau^2T-\tau^3\psi''(1-\tau).
\]
Since $\tau<3/5$, these imply
$\kappa_3<1+6\tau^2T<6v$, or $\kappa_3/(3v)<2$.
Finally,
\[
 r_1=-\frac{\tau P}{b},\qquad
 r_2=r_1-r_1^2-\frac{\tau^2Q}{b}.
\]
Because $r_1<0$,
\[
 r_2-\frac{\kappa_3r_1}{3v}
 <-r_1-r_1^2-\frac{\tau^2Q}{b}
 =\frac{\tau}{b}(P-\tau Q)-r_1^2<0.
\]
Multiplication by $r_1^2/(2v^2)>0$ proves the assertion.
\end{proof}

The proportional theorem settles existence of a small positive sector.
The remaining question is its maximal extent: determine when another
arc or saddle becomes competitive and how that affects coefficient
signs. The cut-density estimates in the half-term theorem normalize by
an exponentially damped coefficient and were proved only for bounded
$m/\sqrt N$; the new coefficient theorem alone does not enlarge that
remainder theorem.
